% =============================================================================
% 057 / 068
% Status: raw
% =============================================================================
% Notes:
%
% Source question:
% 為什麼說元代的漢兒言語並不能僅看作是誕生於金元時代的晚近語言，而是有著更久遠的上游歷史？可否將元代漢兒言語理解為中古漢語隋唐五代西北方音口語的直系後代？
% =============================================================================

\chapter[為什麼說元代的漢兒言語並不能僅看作是誕生於金元時代的晚近語言，而是有著更久遠的上游歷史？可否將元代漢兒言語理解為中古漢語隋唐五代西北方音口語的直系後代？]{為什麼說元代的漢兒言語並不能僅看作是誕生於金元時代的晚近語言，而是有著更久遠的上游歷史？可否將元代漢兒言語理解為中古漢語隋唐五代西北方音口語的直系後代？}
\qaanswer{
首先漢兒言語是中古漢語發展到末期仍然保留大量中古漢語口語語法特徵的語言。而中古漢語口語是西晉滅亡到明初燕王掃北間整整800多年東亞大陸使用的口語。它是鮮卑語，鮮卑語分化出來的北方諸遊牧民族語言和上古漢語不斷融合疊加再融合的產物。中華史觀的人把北亞、東北亞民族按時代割裂成一個個獨立的毫不相干的民族，但事實上自匈奴起，北亞、東北亞的遊牧民族始終是同一波人以及他們彼此融合的子孫。鮮卑當中有很多部落來自早前的匈奴部落，要有一些內亞雜胡部落，他們也不是講統一的語言，與來自匈奴系鮮卑同源的柔然是蒙古的祖先，所以有理由相信他們和在東亞的中原地區建立北魏、北周政權的更早來自匈奴部落的拓拔氏、余文氏所使用的鮮卑語是同源一脈的。可以理解為早一批講匈奴/鮮卑混合語的鮮卑人在漢化中造成了語言的融合，慢慢的成為一種在東亞口頭通用的克里奧爾語。而之後蒙元又在此基礎上二次融入，讓東亞北部的口語更加偏近匈奴/鮮卑一系的語言。值得指出的是在明代建立朱元璋頒布『洪武正正韻』後明政府系統的禁止漢兒言語的使用，特別是朱棣奪權掃北後造成現在的官話區的口語面貌完全看不到漢兒言語留下的痕跡，倒是六南的原生語比如吳語、贛語甚至原生粵語，言從語法、詞彙更偏近中古漢語口語，許多方面能找到漢兒言語的影子。
}

